%! Author = sriadityavedantam
%! Date = 3/25/26

\lesson{14}{Friday 12 September 2025 9:10}{Day 14}

\begin{theorem}[Comparison Test]{
    Given $\abs{a_n}\leq\abs{b_n}$ for all $n$, if $\sum b_n$ converges absolutely, then $\sum a_n$ converges absolutely.
}
    Given $\eps > 0$, there exists $N\in\n$ such that for all $m\geq n\geq N$,
    \[
        \sum_{k=n}^m \abs{b_k} < \eps.
    \]
    This exists by the Cauchy Criterion.
    But
    \[
        \abs{\sum_{k=n}^m a_k}\leq \sum_{k=n}^m\abs{a_k}\leq \sum_{k=n}^m\abs{b_k} < \eps.
    \]
    Hence $\sum a_n$ converges absolutely.
\end{theorem}

\begin{lemma}[Lemma A]{
    If $a < \limsup x_n < b$, then there exists $N\in\n$ such that $n > N$ implies $x_n < b$.
}
    Let $L \coloneqq \limsup x_n$.
    Let $N\in\n$ such that for $n > N$,
    \[
        \abs{y_n - L} < \dfrac{b-L}{2}.
    \]
    Thus
    \[
        y_n < \dfrac{b + L}{2} < b.
    \]
    But if $m\geq n$, then $x_m\leq y_n$.
    Hence for all $m\geq n > N$, we have $x_m < b$.
\end{lemma}

\begin{lemma}[Lemma B]{
    If $a < \limsup x_n$, then for all $N\in\n$ there exists $n > N$ such that $a < x_n$.
}
    Let $L \coloneqq \limsup x_n$.
    Let $N\in\n$.
    Since $y_N \geq L > a$, there exists $m\geq N$ such that $x_m > a$.
    Taking $n \coloneqq m$ proves the claim.
\end{lemma}

\begin{theorem}[Root Test]{
    A series $\sum a_n$ converges absolutely if
    \[
        L \coloneqq \limsup \abs{a_n}^{\dfrac{1}{n}} < 1.
    \]
    It diverges if
    \[
        \limsup \abs{a_n}^{\dfrac{1}{n}} > 1.
    \]
}
    Suppose $L < 1$.
    Let
    \[
        L < r \coloneqq \dfrac{1 + L}{2} < 1.
    \]
    By Lemma A, there exists $N\in\n$ such that $n\geq N$ implies
    \[
        \abs{a_n}^{\dfrac{1}{n}} < r.
    \]
    Hence
    \[
        \abs{a_n} < r^n.
    \]
    Then $\sum \abs{a_n}$ converges by the Comparison Test, since $\sum r^n$ converges.
    Therefore $\sum a_n$ converges absolutely.
    Now suppose
    \[
        \limsup \abs{a_n}^{\dfrac{1}{n}} > 1.
    \]
    Let
    \[
        R \coloneqq \dfrac{1 + \limsup \abs{a_n}^{\dfrac{1}{n}}}{2} > 1.
    \]
    By Lemma B, for all $N\in\n$ there exists $n\geq N$ such that
    \[
        \abs{a_n}^{\dfrac{1}{n}} > R.
    \]
    Thus
    \[
        \abs{a_n} > R^n > 1.
    \]
    Hence $a_n\not\to 0$.
    Therefore $\sum a_n$ diverges.
\end{theorem}

\begin{theorem}[Ratio Test]{
    If
    \[
        \limsup \abs{\dfrac{a_{n+1}}{a_n}} = L < 1,
    \]
    then $\sum a_n$ converges absolutely.
    If
    \[
        \liminf \abs{\dfrac{a_{n+1}}{a_n}} = M > 1,
    \]
    then $\sum a_n$ diverges.
}
    Suppose
    \[
        \limsup \abs{\dfrac{a_{n+1}}{a_n}} = L < 1.
    \]
    Let
    \[
        r \coloneqq \dfrac{1 + L}{2} < 1.
    \]
    By Lemma A, there exists $N\in\n$ such that $n\geq N$ implies
    \[
        \abs{\dfrac{a_{n+1}}{a_n}} < r.
    \]
    Then
    \[
        \abs{a_{N+1}} < r\abs{a_N},
    \]
    and inductively for all $n\geq N$,
    \[
        \abs{a_n}\leq r^{n-N}\abs{a_N}.
    \]
    By the Comparison Test, $\sum \abs{a_N}r^{n-N}$ converges.
    Thus $\sum \abs{a_n}$ converges.
    Hence $\sum a_n$ converges absolutely.
    Now suppose
    \[
        \liminf \abs{\dfrac{a_{n+1}}{a_n}} = M > 1.
    \]
    Let
    \[
        R \coloneqq \dfrac{1 + M}{2} > 1.
    \]
    Then there exists $N\in\n$ such that $n\geq N$ implies
    \[
        \abs{\dfrac{a_{n+1}}{a_n}} > R.
    \]
    Hence
    \[
        \abs{a_{n+1}} > R\abs{a_n}
    \]
    for all $n\geq N$.
    Therefore for all $k\geq N$,
    \[
        \abs{a_k} \geq R^{k-N}\abs{a_N}.
    \]
    Since $R>1$, the sequence $(a_n)$ does not converge to $0$.
    Hence $\sum a_n$ diverges.
\end{theorem}